% Mapa declarativo de retitulado aprobado. No altera los cuerpos científicos.
% Cada entrada procede de una coincidencia estática uno-a-uno y conserva el nivel.
\makeatletter
\DeclareUnicodeCharacter{00B2}{\ensuremath{^{2}}}
\DeclareUnicodeCharacter{00B3}{\ensuremath{^{3}}}
\DeclareUnicodeCharacter{00B7}{\ensuremath{\cdot}}
\DeclareUnicodeCharacter{00D7}{\ensuremath{\times}}
\DeclareUnicodeCharacter{0393}{\ensuremath{\Gamma}}
\DeclareUnicodeCharacter{03B1}{\ensuremath{\alpha}}
\DeclareUnicodeCharacter{03B6}{\ensuremath{\zeta}}
\DeclareUnicodeCharacter{03BE}{\ensuremath{\xi}}
\DeclareUnicodeCharacter{2013}{--}
\DeclareUnicodeCharacter{2074}{\ensuremath{^{4}}}
\DeclareUnicodeCharacter{2076}{\ensuremath{^{6}}}
\DeclareUnicodeCharacter{207B}{\ensuremath{^{-}}}
\DeclareUnicodeCharacter{2080}{\ensuremath{_{0}}}
\DeclareUnicodeCharacter{2081}{\ensuremath{_{1}}}
\DeclareUnicodeCharacter{2082}{\ensuremath{_{2}}}
\DeclareUnicodeCharacter{2083}{\ensuremath{_{3}}}
\DeclareUnicodeCharacter{2085}{\ensuremath{_{5}}}
\DeclareUnicodeCharacter{2087}{\ensuremath{_{7}}}
\DeclareUnicodeCharacter{2088}{\ensuremath{_{8}}}
\DeclareUnicodeCharacter{2089}{\ensuremath{_{9}}}
\DeclareUnicodeCharacter{210F}{\ensuremath{\hbar}}
\DeclareUnicodeCharacter{2192}{\ensuremath{\to}}
\DeclareUnicodeCharacter{21A6}{\ensuremath{\mapsto}}
\DeclareUnicodeCharacter{2212}{\ensuremath{-}}
\DeclareUnicodeCharacter{2218}{\ensuremath{\circ}}
\DeclareUnicodeCharacter{221A}{\ensuremath{\surd}}
\DeclareUnicodeCharacter{2260}{\ensuremath{\ne}}
\DeclareUnicodeCharacter{2261}{\ensuremath{\equiv}}
\DeclareUnicodeCharacter{27E8}{\ensuremath{\langle}}
\DeclareUnicodeCharacter{27E9}{\ensuremath{\rangle}}
\newcommand{\HMTTitleMapEntry}[2]{%
  \edef\HMTTitleProbe{\detokenize{#1}}%
  \ifx\HMTTitleKey\HMTTitleProbe\def\HMTResolvedTitle{#2}\fi}
\newcommand{\HMTResolvePublicTitle}[1]{%
  \def\HMTResolvedTitle{#1}%
  \edef\HMTTitleKey{\detokenize{#1}}%
  % El prefacio fue aprobado como texto autoral literal; su subtítulo no se
  % retitula ni se normaliza.
  \HMTTitleMapEntry{Interacción compleja entre dos máquinas simples}{Interacción compleja entre dos máquinas simples}% ADP102-0001-LITERAL
  \HMTTitleMapEntry{Introducción general}{Arquitectura causal del problema del continuo y alcance matemático-físico de HMT–MD}% ADP102-0002
  \HMTTitleMapEntry{Objeto inicial y orden causal: escala finita, horizonte de observación y publicación arquimediana}{Horizonte proyectivo: dependencia entre escala, frontera y medida}% ADP102-0003
  \HMTTitleMapEntry{Intervalo orientado, marcas, celdas y arista de cierre}{Construcción de la recta finita mediante orden, escala y conservación de la unidad}% ADP102-0004
  \HMTTitleMapEntry{Refinamiento, supervivencia y umbrales de prolongación compatible}{Relación de límite entre infinitésimo, prolongación infinita y umbral proyectivo}% ADP102-0005
  \HMTTitleMapEntry{Evaluación arquimediana de secciones compatibles y reconstrucción de \(\mathbb R\)}{Límite compatible de horizontes finitos y realización de la recta real}% ADP102-0006
  \HMTTitleMapEntry{Núcleo formal de la Holografía Modular Triádica}{Construcción y arquitectura operatoria del núcleo formal APP–TRIT–TPK: soporte aritmético, orientación ternaria y dinámica genealógica}% ADP102-0007
  \HMTTitleMapEntry{Levantamientos y cociclo de acarreo}{Extensiones residuo–cociente y cociclo nonádico de acarreo}% ADP102-0008
  \HMTTitleMapEntry{Bloque central de \(2\) posiciones}{Bloque central bivariado de APP y descomposición en 2 modos espectrales}% ADP102-0009
  \HMTTitleMapEntry{Curvatura mixta de las \(2\) leyes}{Curvatura discreta de la interacción entre las hojas aditiva y multiplicativa}% ADP102-0010
  \HMTTitleMapEntry{Correspondencia entre \(2\) cartas de una misma clase residual}{Conjugación entre las cartas positiva y residual de una misma clase módulo 9}% ADP102-0011
  \HMTTitleMapEntry{Núcleo interior y complemento gnomónico}{Descomposición del ortograma APP en núcleo interior y complemento gnomónico}% ADP102-0012
  \HMTTitleMapEntry{Segmentos marcados, intersticios y célula nonádica}{Construcción de la célula nonádica a partir de segmentos marcados e intersticios}% ADP102-0013
  \HMTTitleMapEntry{Carta posicional finita y proyección residual decimal}{Proyección residual decimal de la carta posicional finita de APP}% ADP102-0014
  \HMTTitleMapEntry{Extracción de la célula de \(9\) fases}{Construcción de la célula nonádica de 9 fases mediante proyección residual}% ADP102-0015
  \HMTTitleMapEntry{Posición, altura y división euclídea}{Coordenadas de posición y altura inducidas por la división euclídea}% ADP102-0016
  \HMTTitleMapEntry{Dominio inaugural de la célula nonádica y exclusión causal de APP, TRIT, TPK y constantes analíticas}{Dominio aritmético autónomo de la célula nonádica y precedencia causal respecto de APP–TRIT–TPK y las constantes analíticas}% ADP102-0017
  \HMTTitleMapEntry{Ley del defecto a módulo fijo y longitud variable}{Invariancia del defecto ciclotómico a módulo fijo bajo variación de longitud}% ADP102-0018
  \HMTTitleMapEntry{Operador ciclotómico primario}{Operador ciclotómico de APP y acción sobre las 12 fibras de tipo A₂}% ADP102-0019
  \HMTTitleMapEntry{Localización del agregado \texorpdfstring{\(54\)}{54}}{Derivación del defecto suma–producto 459−405=54 en el bloque central de APP}% ADP102-0020
  \HMTTitleMapEntry{Resonancia del operador de incidencia en rango \(6\)}{Rango 6 del operador de incidencia y multiplicidad de su modo resonante}% ADP102-0021
  \HMTTitleMapEntry{TRIT como unidad mínima de información topológica: orientación ternaria, memoria de acarreo y orden temporal}{TRIT como unidad ternaria orientada: geometrías internas, orden temporal y condiciones dinámicas de cristal de tiempo}% ADP102-0022
  \HMTTitleMapEntry{Composición y cociclo de acarreo}{Cociclo de acarreo de la composición ternaria orientada}% ADP102-0023
  \HMTTitleMapEntry{Familia universal de álgebras cuadráticas}{Clasificación unificada de las 3 álgebras cuadráticas del TRIT}% ADP102-0024
  \HMTTitleMapEntry{Topological Phase Kernel: base observable y estado enriquecido}{TPK: base observable, estado enriquecido y separación tipada de sus proyecciones}% ADP102-0025
  \HMTTitleMapEntry{La base doble de caminos APP}{Categoría bivariada de caminos APP como base del estado enriquecido TPK}% ADP102-0026
  \HMTTitleMapEntry{Fase nonádica, sector dodecafásico y calendario}{Extensión no escindida C₁₂→C₁₀₈→C₉ y registro dodecafásico de la fase nonádica}% ADP102-0027
  \HMTTitleMapEntry{Cilindro y frontera}{Compatibilidad cilíndrica y firma de frontera del estado enriquecido}% ADP102-0028
  \HMTTitleMapEntry{Transporte de fibras}{Operador cardinal \(T_d\) de transporte entre fibras enriquecidas con conservación de ruta y memoria}% ADP102-0029
  \HMTTitleMapEntry{Relaciones de extensión}{Relación tipada \(\operatorname{Ext}_r\) de prolongación entre fibras enriquecidas}% ADP102-0030
  \HMTTitleMapEntry{Proyecciones y pérdidas tipadas}{Cocientes observables del estado enriquecido y campos genealógicos conservados}% ADP102-0031
  \HMTTitleMapEntry{Árbol operatorio del contenedor}{Inserción de las 6 familias fibradas del estado enriquecido en la jerarquía operatoria del TPK}% ADP102-0032
  \HMTTitleMapEntry{Álgebra de operadores del Topological Phase Kernel y dinámica causal de actualización del estado enriquecido}{Jerarquía tipada de las 34 construcciones canónicas heterogéneas y dinámica causal de actualización del TPK}% ADP102-0033
  \HMTTitleMapEntry{Clasificación del estado enriquecido en \(8\) clases estructurales invariantes}{Jerarquía tipada L₀–L₇ de las 34 construcciones canónicas heterogéneas del TPK}% ADP102-0034
  \HMTTitleMapEntry{Descomposición funcional en \(14\) agrupaciones compatibles con los operadores APP–TRIT–TPK}{Agrupación funcional de las 34 construcciones canónicas en 14 familias de composición APP–TRIT–TPK}% ADP102-0035
  \HMTTitleMapEntry{Dominio y tipado de las \(34\) entradas del operador de actualización}{Clases matemáticas, dominios, codominios y dependencias causales de las 34 construcciones canónicas}% ADP102-0036
  \HMTTitleMapEntry{Composición causal de la actualización del estado enriquecido}{Actualización del estado enriquecido mediante \(U_t=\operatorname{Upd}_t\circ\operatorname{Tra}_t\circ\operatorname{Sel}_t\)}% ADP102-0037
  \HMTTitleMapEntry{Separación tipada entre operadores de transporte, memoria, calendario, extracción y retorno}{Separación tipada entre operadores, estados, proyecciones y publicaciones del TPK}% ADP102-0038
  \HMTTitleMapEntry{Descomposición operatorial de los canales de incidencia \texorpdfstring{\(90/120\)}{90/120}: selector temporal, generador nilpotente y componente angular}{Incidencia trítica 90/120 y separación tipada de los operadores temporal, nilpotente y angular}% ADP102-0039
  \HMTTitleMapEntry{Distribución de los índices \(90/120\) entre los \(6\) objetos del sistema de incidencia}{Residencias tipadas de los índices 90/120 en 6 objetos operatorios no equivalentes}% ADP102-0040
  \HMTTitleMapEntry{Defecto normalizado de la incidencia}{Defecto normalizado 1−90/120=1/4 de la incidencia trítica}% ADP102-0041
  \HMTTitleMapEntry{Invariantes de tipado}{Invariantes de dominio, orientación y memoria de las realizaciones 90/120}% ADP102-0042
  \HMTTitleMapEntry{Estado local y realimentación de hoja}{Operador local de transición orientada y actualización de la hoja APP}% ADP102-0043
  \HMTTitleMapEntry{Cilindros de los \(2\) lectores}{Sistemas cilíndricos de los 2 lectores aritméticos y compatibilidad de truncación}% ADP102-0044
  \HMTTitleMapEntry{Descomposición del Topological Phase Kernel en operadores de avance, giro, acarreo y memoria}{Factorización de la transición TPK como \(U_t=\operatorname{Upd}_t\circ\operatorname{Tra}_t\circ\operatorname{Sel}_t\)}% ADP102-0045
  \HMTTitleMapEntry{Soporte ordenado y coordenada cíclica}{Coordenada cíclica sobre el soporte temporal orientado de 12 sectores}% ADP102-0046
  \HMTTitleMapEntry{Extensión cíclica no escindida}{Extensión no escindida 0→C₁₂→C₁₀₈→C₉→0 y cociclo de acarreo}% ADP102-0047
  \HMTTitleMapEntry{Extractor firmado de eventos}{Extracción del estado dodecafásico firmado a partir del registro temporal enriquecido}% ADP102-0048
  \HMTTitleMapEntry{\(2\) palabras dodecafásicas de distinta procedencia}{Distinción genealógica entre 2 palabras dodecafásicas de igual período}% ADP102-0049
  \HMTTitleMapEntry{Refinamiento mínimo por microinstantes}{Subdivisión temporal mínima compatible con la semiconjugación}% ADP102-0050
  \HMTTitleMapEntry{\(3\) operaciones cuaternarias no intercambiables}{Conmutadores de 3 operaciones cuaternarias y separación de sus acciones}% ADP102-0051
  \HMTTitleMapEntry{Lector trítico del estado orientado: orden temporal y transporte de acarreo}{Lector trítico orientado con transporte de cociente y acarreo}% ADP102-0052
  \HMTTitleMapEntry{Cociente mínimo de memoria predictiva}{Cociente de memoria de orden 1 y suficiencia predictiva de la proyección modal}% ADP102-0053
  \HMTTitleMapEntry{Memoria vertical no acotada}{Crecimiento no acotado de la profundidad de memoria bajo retornos de fase}% ADP102-0054
  \HMTTitleMapEntry{Realizaciones geométricas del TRIT: regímenes elíptico, parabólico e hiperbólico y doble orientación temporal}{TRIT como estado ternario orientado: memoria de acarreo, orden temporal y 3 regímenes cuadráticos}% ADP102-0055
  \HMTTitleMapEntry{APP, TRIT y TPK: soporte, orientación y transporte}{Composición APP–TRIT–TPK: soporte aritmético, orientación ternaria y transporte con memoria}% ADP102-0056
  \HMTTitleMapEntry{TPK como sistema de transporte}{TPK como dinámica de selección, transporte, memoria, retorno y despliegue}% ADP102-0057
  \HMTTitleMapEntry{Periodicidad temporal de (108) fases y leyes de retorno a (27), (54) y (108) iteraciones}{Periodicidad temporal de 108 fases y retornos a 27, 54 y 108 iteraciones}% ADP102-0058
  \HMTTitleMapEntry{APP, TRIT y TPK: definiciones rectoras}{Composición generativa APP→TRIT→TPK y construcción del estado enriquecido}% ADP102-0059
  \HMTTitleMapEntry{Consecuencias y publicaciones del continuo construido}{Conexión nonádica conjunta Γ₉ como correspondencia compatible del sistema inverso}% ADP102-0060
  \HMTTitleMapEntry{Condicionamiento como poda de futuros}{Condicionamiento de la medida mediante restricción del árbol de prolongaciones futuras}% ADP102-0061
  \HMTTitleMapEntry{Holonomía nonádica \texorpdfstring{\(\Gamma_9\)}{Γ₉} del Topological Phase Kernel: fibra \texorpdfstring{\(\mathbb F_3^6\)}{F₃⁶}, prolongación coinductiva y retorno de fase con memoria de estado}{Holonomía nonádica Γ₉ del Topological Phase Kernel: transporte del estado enriquecido, prolongación coinductiva y retorno 9→1 con avance de memoria}% ADP102-0062
  \HMTTitleMapEntry{Conexión nonádica conjunta, compresión y memoria terminal}{Holonomía conjunta Γ₉: compresión observable T, componente ortogonal η y conservación terminal de memoria}% ADP102-0063
  \HMTTitleMapEntry{Relaciones locales y holonomía}{Relaciones locales de transición y construcción de la holonomía nonádica Γ₉}% ADP102-0064
  \HMTTitleMapEntry{Obstrucciones de la holonomía \(\Gamma_9\): retorno \(9\to1\), índices de composición y conservación de fase, acarreo, frontera, ruta y memoria}{Obstrucciones de la holonomía Γ₉: retorno 9→1, índices de composición y conservación de fase, acarreo, frontera, ruta y memoria}% ADP102-0065
  \HMTTitleMapEntry{Descomposición de la microfibra de clausura de \texorpdfstring{\(\pi\)}{π} en \texorpdfstring{\(5\)}{5} regiones orientadas: multiplicidad \texorpdfstring{\(432+4\cdot144=1008\)}{432+4·144=1008}, eje \texorpdfstring{\(\varphi\)}{φ} e involución especular \texorpdfstring{\(\pi/e\)}{π/e}}{Clasificación orientada de la microfibra de clausura de π: 5 regiones, multiplicidad 432+4·144=1008, eje φ e involución π/e}% ADP102-0066
  \HMTTitleMapEntry{Generación coinductiva de \texorpdfstring{\(\pi\)}{π}, \texorpdfstring{\(e\)}{e} y \texorpdfstring{\(\varphi\)}{φ} mediante refinamiento nonádico compatible}{Generación coinductiva de π, e y φ mediante refinamiento nonádico compatible}% ADP102-0067
  \HMTTitleMapEntry{Elevación visible y cadena de bloques}{Prolongación coinductiva de la emisión exponencial mediante bloques compatibles}% ADP102-0068
  \HMTTitleMapEntry{Firma decimal y cuadrado posicional}{Construcción de la firma decimal mediante el cuadrado posicional ternario–decimal}% ADP102-0069
  \HMTTitleMapEntry{Retorno visible y ruptura de memoria}{Retorno de fase con actualización no trivial de la memoria del estado}% ADP102-0070
  \HMTTitleMapEntry{Certificado adversarial de la generación de \(e\): unicidad orbital, recurrencia factorial, horquillas racionales y prolongación cilíndrica}{Unicidad orbital, recurrencia factorial y convergencia cilíndrica de la generación de e}% ADP102-0071
  \HMTTitleMapEntry{Elevación por bloques y 1.ª frontera coinductiva}{Prolongación por bloques y construcción de la 1.ª frontera coinductiva}% ADP102-0072
  \HMTTitleMapEntry{Del calendario trítico al multigrafo de incidencia}{Construcción del multigrafo de incidencia a partir del calendario trítico}% ADP102-0073
  \HMTTitleMapEntry{Caracterización, horquillas y prolongación}{Caracterización de la autoescala mediante horquillas racionales y prolongación coinductiva}% ADP102-0075
  \HMTTitleMapEntry{Selección diagonal mínima}{Selección diagonal mínima de cilindros compatibles y convergencia a precisión arbitraria}% ADP102-0076
  \HMTTitleMapEntry{Transferencia Markov compatible con todo el límite inverso}{Dinámica de Markov como cociente modal de memoria de orden 1 de la holonomía Γ₉, compatible con el límite inverso}% ADP102-0077
  \HMTTitleMapEntry{2.º sistema de coordenadas: diferencias de pasos \(3\) y \(4\)}{Sistema diferencial de coordenadas inducido por los operadores de paso 3 y 4}% ADP102-0078
  \HMTTitleMapEntry{Obstrucciones de la prolongación afín de \(\alpha\): telescopía, inversión de \(K\), frontera \(662/663\) y aislamiento de datos objetivo}{Obstrucciones de la prolongación afín de α: telescopía, inversión de K, frontera 662/663 y aislamiento de datos objetivo}% ADP102-0079
  \HMTTitleMapEntry{Construcción de \texorpdfstring{\(\sqrt{-1}\)}{√(-1)} mediante \texorpdfstring{\(J^2=-I\)}{J²=-I} y coordinación de las geometrías elíptica, parabólica e hiperbólica}{Construcción de √(−1) y familia exponencial trítica de Euler: operadores J_τ y coordinación de los regímenes elíptico, parabólico e hiperbólico}% ADP102-0080
  \HMTTitleMapEntry{El Núcleo de Transducción Hensel--Witt}{Núcleo de transducción Hensel–Witt: dominio, operación y preservación genealógica}% ADP102-0081
  \HMTTitleMapEntry{Núcleo de Transducción Hensel–Witt: microregiones de \(\pi\), holonomía \(\Gamma_9\), registro dodecafásico y transporte hexada–octada}{Núcleo de Transducción Hensel–Witt: microregiones de π, holonomía Γ₉, registro dodecafásico y transporte hexada–octada}% ADP102-0082
  \HMTTitleMapEntry{Del código ternario al sistema de Witt}{Transducción del código ternario al sistema de Witt}% ADP102-0083
  \HMTTitleMapEntry{Hexada, octada y el corredor \texorpdfstring{\(90/120\)}{90/120}}{Inclusión hexada–octada y transporte de incidencia por los canales 90/120}% ADP102-0084
  \HMTTitleMapEntry{Extractor de escala triádico}{Operador triádico de extracción de escala y su imagen racional}% ADP102-0085
  \HMTTitleMapEntry{Defecto de potencias primas y operador número}{Representación del defecto de potencias primas mediante el operador número}% ADP102-0086
  \HMTTitleMapEntry{Ablaciones del selector y de la continuación}{Condiciones de unicidad del selector regional y de la prolongación analítica}% ADP102-0087
  \HMTTitleMapEntry{Cadena de generación \(\mathcal F_\pi\to b_\pi\to\rho_\pi\to D_A\to C_*^{\rm HMT}\to\Gamma_{6\to8}\): selección regional, prolongación determinantal y coordenada orientada}{Cadena de generación \(\mathcal F_\pi\to b_\pi\to\rho_\pi\to D_A\to C_*^{\mathrm{HMT}}\to\Gamma_{6\to8}\): selección regional, prolongación determinantal y coordenada orientada}% ADP102-0088
  \HMTTitleMapEntry{Descomposición incidencial hexada--octada en los sectores \texorpdfstring{\(q_{90}\)}{q₉₀} y \texorpdfstring{\(q_{120}\)}{q₁₂₀}: entrelazador \texorpdfstring{\(6\to8\)}{6→8} y estructura precursora del espectro de área}{Canales de incidencia 90/120 del Topological Phase Kernel: semicoronas tríticas, factorización de hojas y realización hexada–octada del morfismo 6→8}% ADP102-0089
  \HMTTitleMapEntry{Geometría dodecafásica de mezcla: antecedente angular CKM--CP y lector precursor del electrón}{Operador de mezcla dodecafásico: fase CKM–CP y construcción del lector predimensional del electrón}% ADP102-0090
  \HMTTitleMapEntry{Morfismos arquimediano y excepcional desde un dominio enriquecido común}{Principio de Ambivalencia: involución de los lectores areal y volumétrico sobre el estado enriquecido}% ADP102-0091
  \HMTTitleMapEntry{La operación local sobre las 729 palabras}{Dominio enriquecido común y codominios inequivalentes de las 2 proyecciones}% ADP102-0092
  \HMTTitleMapEntry{Naturalidad y compatibilidad a toda profundidad}{Separación conjunta de historias por las proyecciones arquimediana y excepcional}% ADP102-0093
  \HMTTitleMapEntry{Acoplamiento postgenerativo con el cierre operatorial de Euler}{Invariante de incidencia común entre el estado enriquecido y el sistema de Witt}% ADP102-0094
  \HMTTitleMapEntry{Globalización natural de las proyecciones arquimediana y excepcional sobre un sistema inverso común}{Descomposición de la microfibra de π en 5 regiones orientadas: 4 componentes periféricas y 1 eje fijo}% ADP102-0095
  \HMTTitleMapEntry{Sistemas compatibles a profundidad finita}{Prolongación reticular de la incidencia y construcción de pantallas dimensionales}% ADP102-0096
  \HMTTitleMapEntry{Calibre de la publicación excepcional}{Factorización del lector de incidencia a través del estado enriquecido}% ADP102-0097
  \HMTTitleMapEntry{Naturalidad conjunta de las publicaciones del estado genealógico}{Doble publicación del estado enriquecido: evaluación arquimediana y conservación de incidencia}% ADP102-0098
  \HMTTitleMapEntry{Separación conjunta de valor e incidencia excepcional}{Unicidad del calibre inicial de la proyección excepcional}% ADP102-0099
  \HMTTitleMapEntry{Antecedente no conmutativo común de la doble proyección}{Álgebra no conmutativa común y factorización de las 2 proyecciones}% ADP102-0100
  \HMTTitleMapEntry{Del espacio de historias a la recta real}{Cociente del espacio de historias compatibles y realización de la recta real}% ADP102-0101
  \HMTTitleMapEntry{Media arquimediana y doble proyección del mismo estado}{Compatibilidad de la evaluación arquimediana media con la doble proyección del estado enriquecido}% ADP102-0102
  \HMTTitleMapEntry{Calendario unitario, memoria modular y modo estable}{Modo estable del calendario unitario bajo actualización de memoria modular}% ADP102-0103
  \HMTTitleMapEntry{Publicaciones analíticas de historias HMT}{Realizaciones analíticas de historias HMT mediante las cartas arquimediana y bidireccional}% ADP102-0104
  \HMTTitleMapEntry{Consecuencia para una acción positiva}{Positividad de la acción inducida por el entrelazador dodecafásico–Witt}% ADP102-0105
  \HMTTitleMapEntry{Del salto espectral \texorpdfstring{\(4\)}{4} al retículo de Leech: despliegue APP \texorpdfstring{\(4\to2\to6\to54\)}{4→2→6→54} y corredor Paley--Witt--Mathieu--Golay--Niemeier}{Del salto espectral 4 al retículo de Leech: despliegue APP 4→2→6→54 y corredor Paley–Witt–Mathieu–Golay–Niemeier}% ADP102-0106
  \HMTTitleMapEntry{Naturaleza de la estructura obtenida}{Clasificación algebraica de la estructura ternaria autodual obtenida}% ADP102-0107
  \HMTTitleMapEntry{Ley de incidencia \texorpdfstring{\(4+1\)}{4+1} sobre una tétrada}{Descomposición 4+1 de la incidencia de una tétrada en componentes periféricas y axial}% ADP102-0108
  \HMTTitleMapEntry{Módulo y forma discriminantes}{Descomposición de la incidencia de una tétrada en 4 componentes periféricas y 1 componente axial}% ADP102-0109
  \HMTTitleMapEntry{\(12\) factores y aplicación de pegado}{Correspondencia entre estratos de incidencia y órbita libre del triple positivo}% ADP102-0110
  \HMTTitleMapEntry{Carácter y núcleo de congruencia}{Pegado discriminante de 12 factores A₂ y construcción de la extensión integral}% ADP102-0111
  \HMTTitleMapEntry{Curvaturas mixtas de APP}{Sistema ordenado de incidencias y orientación de la cámara radial}% ADP102-0112
  \HMTTitleMapEntry{Representación de Weyl--Heisenberg y fase central}{Construcción del carácter radial y determinación de su núcleo de congruencia}% ADP102-0113
  \HMTTitleMapEntry{Estatuto exacto de la jerarquía \(4\to2\to6\to54\) y de la identificación reticular.}{Estatuto exacto de la jerarquía 4→2→6→54 y de la identificación reticular.}% ADP102-0114
  \HMTTitleMapEntry{Traza graduada del generador ternario}{Traza graduada de la acción del generador ternario sobre \(V^{\natural}\)}% ADP102-0115
  \HMTTitleMapEntry{Palabra de soporte total e isometría fijo-libre}{Isometría sin puntos fijos inducida por la palabra ternaria de soporte total}% ADP102-0116
  \HMTTitleMapEntry{Construcción del espacio de Hilbert a partir de la conexión nonádica y sus operadores de fase}{Realización proyectiva de Bloch a partir de la estructura compleja interna}% ADP102-0117
  \HMTTitleMapEntry{La rama marcada cambia de tipo}{Inmersión por esquinas de una fase marcada: ideal compacto \(\mathcal K(\mathcal H_\infty)\) y factor \(\mathcal B(\mathcal H_\infty)\) de tipo \(\mathrm I_\infty\)}% ADP102-0118
  \HMTTitleMapEntry{Residuo visible y cociente de frontera}{Extensión operatoria de APP: modos visibles y modos de memoria}% ADP102-0119
  \HMTTitleMapEntry{Involución de dualidad T}{Completación 1000=729+243+27+1 y estratificación de la unidad}% ADP102-0120
  \HMTTitleMapEntry{Completación estratificada de la unidad}{Descomposición 243=1+22+90+120+10 y tipado de sus componentes}% ADP102-0121
  \HMTTitleMapEntry{El bloque \texorpdfstring{\(243\)}{243} como frontera y bola ternaria perfecta}{Derivación de los canales de constantes a partir de la descomposición del bloque 243}% ADP102-0122
  \HMTTitleMapEntry{Descomposición interna del bloque \texorpdfstring{\(243\)}{243}}{Construcción de la dimensión 11 mediante lectores residuales compatibles}% ADP102-0123
  \HMTTitleMapEntry{Dimensión \texorpdfstring{\(11\)}{11} y lecturas compatibles}{Clausura del registro genealógico bajo la involución de dualidad T}% ADP102-0124
  \HMTTitleMapEntry{Identidades verificables de la construcción}{Identidades exactas de involución, modularidad y conservación residual}% ADP102-0125
  \HMTTitleMapEntry{Nivel, intercambio lineal de masas y modularidad: \(3\) controles de no fusión}{Separación por nivel, intercambio lineal de masas y modularidad de las 3 realizaciones}% ADP102-0126
  \HMTTitleMapEntry{\(2\) descendientes del código ternario extendido}{Reducción dimensional 26→24→11→10 y coordinación con las pantallas de teoría M}% ADP102-0127
  \HMTTitleMapEntry{La representación icosaédrica de dimensión \(12\)}{Coordinación de las pantallas exactas con la realización física de teoría M}% ADP102-0128
  \HMTTitleMapEntry{Reducción isotrópica de rango \(26\to24\)}{Representación operatoria de APP y terna espectral supersimétrica en dimensión 11}% ADP102-0129
  \HMTTitleMapEntry{Coordinación del corredor Moonshine con las pantallas de teoría M: gradación, retículo lorentziano y doble realización del estado HMT}{Teorema de compatibilidad entre el corredor Moonshine y las pantallas de teoría M mediante gradación y transporte reticular}% ADP102-0130
  \HMTTitleMapEntry{Antecedente común de lectores: refinamiento espectral y límite de dilatación conjunta}{Refinamiento espectral común y límite de la dilatación conjunta de los lectores matriciales}% ADP102-0131
  \HMTTitleMapEntry{Subestructuras, objeto global y proyecciones}{Construcción de 4 cuadrados latinos cuánticos cíclicos en una clase operatoria elemental}% ADP102-0132
  \HMTTitleMapEntry{Correspondencias directas e inversas entre los formalismos}{Realización de 3 contextos operatorios en un cuadrado cuántico (9,3)}% ADP102-0133
  \HMTTitleMapEntry{Extensión de la bigraduación a 3 índices}{Representación del atlas genealógico completo en un cuadrado cuántico (12,3)}% ADP102-0134
  \HMTTitleMapEntry{4 cuadrados latinos cuánticos cíclicos de clase operatorial fácil}{Clasificación cíclica de 4 cuadrados latinos cuánticos en una clase operatoria elemental}% ADP102-0135
  \HMTTitleMapEntry{Lema cíclico de cuadrados mágicos cuánticos asociados a una medida operatorial positiva}{Bigraduación de la familia genealógico–matricial}% ADP102-0136
  \HMTTitleMapEntry{Conos matriciales, rangos y extensiones completamente positivas}{Representación matricial de la composición genealógica APP–TRIT–TPK}% ADP102-0137
  \HMTTitleMapEntry{Construcción explícita de cubos mágicos cuánticos y compatibilidad de sus márgenes}{TPK como categoría de rutas con selección, transporte, memoria y retorno}% ADP102-0138
  \HMTTitleMapEntry{Combinaciones mágicas de mapas CP}{Composiciones completamente positivas compatibles con la bigraduación genealógico–matricial}% ADP102-0139
  \HMTTitleMapEntry{Jerarquía espectral \(4\to2\to6\to54\) del bloque central de APP}{Realización metrológica posterior de la doble proyección del estado enriquecido}% ADP102-0140
  \HMTTitleMapEntry{Refinamiento nonádico mediante esperanza condicional}{Descomposición simpléctica 12→4×3 de las coordenadas de fase}% ADP102-0141
  \HMTTitleMapEntry{Taxonomía funcional de \(114\) expresiones numéricas y \(106\) lecturas operatoriales}{Clasificación operatoria de 114 expresiones numéricas por dominio, generador, invariante y codominio}% ADP102-0142
  \HMTTitleMapEntry{Transferencia de fase sobre el multigrafo de bloques completos}{Microfibra de clausura de π: 5 regiones orientadas, multiplicidades \(432+4\cdot144=1008\) y proyección común \(w_\pi\)}% ADP102-0143
  \HMTTitleMapEntry{\(5\) ciclos orientados de clausura de la microfibra}{Publicación de memoria reducida de un bloque nonádico completo mediante el multigrafo de transferencia de fase}% ADP102-0144
  \HMTTitleMapEntry{Operador de incidencia y media vuelta orientada}{Holonomías orientadas de clausura de las 5 regiones de π}% ADP102-0145
  \HMTTitleMapEntry{Composición racional y compensador \texorpdfstring{\(239\)}{239}}{Determinación de la pendiente 1/5 por la multiplicidad de la microfibra regional de π}% ADP102-0146
  \HMTTitleMapEntry{Identidad gaussiana y \(1^{\mathrm a}\) media vuelta orientada}{Composición racional de la media vuelta y determinación del término compensador 239}% ADP102-0147
  \HMTTitleMapEntry{Prolongación coinductiva de la sección de \(\pi\) mediante el sistema inverso de historias supervivientes}{Filtración de F₃⁶ por admisibilidad regional y selección de subcilindros compatibles}% ADP102-0148
  \HMTTitleMapEntry{Sistema inverso de historias compatibles}{Unicidad de la sección diagonal mínima en el sistema de horquillas racionales}% ADP102-0149
  \HMTTitleMapEntry{Firma posicional decimal y cuadrado de compatibilidad}{Prolongación de la sección exponencial mediante la cadena de bloques visibles}% ADP102-0150
  \HMTTitleMapEntry{Retorno de fase con cambio de estado y conservación de la memoria recorrida}{Construcción de la firma decimal y del cuadrado posicional de la sección exponencial}% ADP102-0151
  \HMTTitleMapEntry{Semigrupo factorial del carácter exponencial: recurrencia \(a_n=1/n!\)}{Retorno de fase con desigualdad de estados enriquecidos: \(g(K+9)=g(K)\) y \(B_{K+9}\ne B_K\)}% ADP102-0152
  \HMTTitleMapEntry{Certificado interno de unicidad orbital, recurrencia factorial, convergencia y prolongación cilíndrica}{Prolongación coinductiva de e mediante la holonomía nonádica Γ₉ del TPK}% ADP102-0153
  \HMTTitleMapEntry{Rayo positivo invariante y unicidad de la autoescala}{Construcción del multigrafo de incidencia modal a partir del calendario trítico}% ADP102-0154
  \HMTTitleMapEntry{Certificado interno de unicidad, compatibilidad y convergencia de la autoescala de Perron}{Caracterización de φ por el rayo de Perron, encaje de horquillas y prolongación coinductiva}% ADP102-0156
  \HMTTitleMapEntry{Derivaciones internas de la constante de estructura fina \(\alpha\): cierre entero dodecafásico reversible y raíz analítica única del núcleo \(E_{21/22}\)}{Clausura dodecafásica de \(\alpha_{\mathrm{HMT}}\) y caracterización analítica mediante \(E_{21/22}\): generación reversible, raíz simple y prolongación nonádica de orden 9}% ADP102-0157
  \HMTTitleMapEntry{Estado terminal común y 2 lecturas internas coordinadas}{Evaluaciones terminal y dodecafásica del estado enriquecido común: bloques \(B_{\mathrm{term}}\) y carta firmada \(U_{12}^{\mathrm{sgn}}\)}% ADP102-0158
  \HMTTitleMapEntry{Vía A: cierre entero dodecafásico}{Generación reversible de \(\alpha_{\mathrm{HMT}}\) mediante la clausura dodecafásica entera}% ADP102-0159
  \HMTTitleMapEntry{Dominio Hensel--Witt, mapas de Teichmüller y objetivo racional \(-1/12\)}{Operador Hensel–Witt de transducción entre microregiones de π, holonomía Γ₉, registro dodecafásico e incidencia hexada–octada}% ADP102-0160
  \HMTTitleMapEntry{Aplicación de Teichmüller del código ternario al anillo de Witt}{Construcción del sistema de Witt desde la codificación ternaria compatible}% ADP102-0161
  \HMTTitleMapEntry{Incidencia hexada--octada y canales modulares \(90/120\)}{Realización hexada–octada de los canales 90/120 previamente derivados de la incidencia trítica del TPK}% ADP102-0162
  \HMTTitleMapEntry{4 rutas compatibles hacia \texorpdfstring{\(-1/12\)}{-1/12}}{4 realizaciones operatorias de −1/12 y separación de sus dominios}% ADP102-0163
  \HMTTitleMapEntry{Derivación determinantal de la escala de acción: forma normal de Smith, conúcleo de orden \(16\) y valoración \(10^{-34}\)}{Derivación determinantal de la escala absoluta de acción: forma normal de Smith, conúcleo de orden 16 y valoración 10⁻³⁴}% ADP102-0164
  \HMTTitleMapEntry{Medida de cociente y carácter positivo de la acción: hilbertización canónica, defecto regional--dodecafásico y ley functorial sobre caminos}{Construcción de la medida positiva de acción sobre el cociente de emisiones: hilbertización, defecto regional–dodecafásico y functorialidad de caminos}% ADP102-0165
  \HMTTitleMapEntry{Derivación autodual de la constante de Planck: periodicidad temporal de \(108\) fases, acción y longitud de Compton}{Derivación autodual de la constante de Planck y de su recta de acción: periodicidad temporal de 108 fases, cierre de Compton y covariancia bidireccional}% ADP102-0166
  \HMTTitleMapEntry{Equivalencia de normalizaciones sobre la identidad autodual}{Conmutatividad de las normalizaciones angular, temporal y gravitatoria de la identidad autodual}% ADP102-0167
  \HMTTitleMapEntry{Operador angular predimensional: descomposición en \(27\) sectores, radio espectral y composición efectiva de la doble proyección}{Operador de fase anterior a la transducción dimensional: descomposición en 27 sectores, radio espectral y factorización de la doble proyección}% ADP102-0168
  \HMTTitleMapEntry{Caracteres cuadráticos y constantes espectrales: Catalán, Apéry y torre de Euler--Stieltjes}{Caracteres cuadráticos y constantes espectrales: Catalán, Apéry, Euler–Mascheroni y jerarquía de Stieltjes}% ADP102-0169
  \HMTTitleMapEntry{Criticidad unimodal de la familia dodecafásica: perfil analítico exacto, aproximantes de Feigenbaum y reducción de la universalidad al certificado hiperbólico}{Criticidad unimodal de la familia dodecafásica: perfil analítico exacto, aproximantes de Feigenbaum y certificado hiperbólico de renormalización}% ADP102-0170
  \HMTTitleMapEntry{Aritmética del modo \(30\): identidades de Rogers--Ramanujan, normalizador \(899\) y levantamiento nonádico de Pell}{Aritmética de la familia modular dodecafásica de periodo 30: identidades de Rogers–Ramanujan, normalizador 899 y levantamiento nonádico de Pell}% ADP102-0171
  \HMTTitleMapEntry{Teorema rector de la red radical y la ley potencia--vacancia}{Clasificación de la red radical 3→6→9 y derivación de la ley potencia–vacancia}% ADP102-0172
  \HMTTitleMapEntry{Órbita cíclica \(369\to693\to936\to369\)}{Órbita cíclica 369→693→936→369}% ADP102-0173
  \HMTTitleMapEntry{Núcleo \(7227\), descomposición espectral y enlace con la cruz radical}{Determinación del núcleo aritmético 7227 por incidencia con la cruz radical de APP}% ADP102-0174
  \HMTTitleMapEntry{Codificación mecánica de las vacancias por las potencias nonádicas}{Codificación sturmiana de las vacancias inducida por las potencias nonádicas}% ADP102-0175
  \HMTTitleMapEntry{Raíz interna de \(-1\), estructura exponencial y geometrías elíptica, parabólica e hiperbólica}{Familia exponencial triádica \(J_\tau^2=-\tau I\): realizaciones elíptica, parabólica e hiperbólica de la identidad de Euler}% ADP102-0176
  \HMTTitleMapEntry{Realización continua del tipo cuadrático inducido}{Extensión continua del operador cuadrático común y conservación de su tipo espectral}% ADP102-0177
  \HMTTitleMapEntry{Soporte operatorio común de \(3\) lecturas geométricas}{Conjugación de las realizaciones aritmética, angular y torsional mediante el operador cuadrático común}% ADP102-0178
  \HMTTitleMapEntry{Cadena constitutiva del sector gravitatorio: escala cronológica \(108\), reducción tensorial, holonomía y evaluación metrológica de \(G\)}{Composición causal de los invariantes de fase, acción y longitud en la derivación de G}% ADP102-0179
  \HMTTitleMapEntry{Cardinalidades de incidencia 90 y 120}{Derivación trítica de las cardinalidades 90/120 y realización posterior en la incidencia hexada–octada}% ADP102-0180
  \HMTTitleMapEntry{Separación causal entre el espectro de área, la memoria modular y la realización Cartan--Holst}{Distinción tipada entre Barbero–área, área–Hamiltoniano modular y realización Cartan–Holst}% ADP102-0182
  \HMTTitleMapEntry{Transductor de Hensel no filtrado y desplazamiento completo sobre las fibras ternarias}{Transductor integral de Hensel como operador de desplazamiento completo}% ADP102-0183
  \HMTTitleMapEntry{Sobreyectividad de la carta arquimediana no filtrada sobre \(\mathbb R\)}{Cobertura exacta de la carta real antecedente mediante secciones compatibles}% ADP102-0184
  \HMTTitleMapEntry{2 ejes de clasificación: valor visible y genealogía de transporte}{Clasificación bifactorial del número por valor arquimediano y genealogía de transporte}% ADP102-0185
  \HMTTitleMapEntry{Posición de \texorpdfstring{\(\pi,e,\varphi\) y \(\alpha\)}{pi, e, phi y alpha} en la clasificación}{Posición de π,e,φ y α en la clasificación}% ADP102-0186
  \HMTTitleMapEntry{Carta de Hensel no filtrada, supervivencia y frontera global}{Carta antecedente, supervivencia coinductiva y construcción de la frontera global}% ADP102-0187
  \HMTTitleMapEntry{La familia operatorial de Euler como gramática de cierre}{Factorización del cierre genealógico del número mediante el sistema operatorio de Euler}% ADP102-0188
  \HMTTitleMapEntry{Dependencias matemáticas.}{Dependencia causal entre sección inversa, evaluación arquimediana y conservación excepcional de la incidencia}% ADP102-0189
  \HMTTitleMapEntry{\(4\) exigencias de elección en el sistema proyectivo}{Condiciones de elección necesarias para la prolongación global de secciones compatibles}% ADP102-0190
  \HMTTitleMapEntry{Decidibilidad efectiva de cada ventana finita de multiplicidades proyectivas}{Decidibilidad algorítmica de cada nivel finito del sistema de prefijos}% ADP102-0191
  \HMTTitleMapEntry{Relación exacta con las \(9\) fases del Topological Phase Kernel}{Restricción de la obstrucción uniforme al sistema temporal nonádico de 9 fases}% ADP102-0192
  \HMTTitleMapEntry{De la consistencia constructible a la independencia por forcing}{Consistencia relativa en L e independencia mediante forzamiento}% ADP102-0193
  \HMTTitleMapEntry{Optimización minimax sobre ventanas finitas de prefijos compatibles}{Teorema minimax en cada nivel finito del poset de prefijos compatibles}% ADP102-0194
  \HMTTitleMapEntry{La cadena que el puente recibe demostrada}{Antecedentes teoremáticos del funtor de transducción genealógico–metrológica}% ADP102-0195
  \HMTTitleMapEntry{Derivación bidireccional del exponente electrónico desde la ruta enriquecida}{Banda de realización del electrón y separación entre masa basal y clausura beta}% ADP102-0196
  \HMTTitleMapEntry{\(2\) ramas causales y su confluencia}{Confluencia de las ramas determinantal y dodecafásica de la escala de acción}% ADP102-0197
  \HMTTitleMapEntry{Carta metrológica posterior y evaluación dimensional de los invariantes generados}{Sistemas de transducción entre número genealógico y magnitud física}% ADP102-0198
  \HMTTitleMapEntry{Principio de Ambivalencia: involución de los lectores areal y volumétrico, equivalencia local y orientación temporal conjugada}{Principio de Ambivalencia en Mecánica Dimensional: involución areal–volumétrica, sistemas autoinerciales y orientación temporal conjugada}% ADP102-0199
  \HMTTitleMapEntry{Condiciones de una evaluación prospectiva}{Condiciones de evaluación prospectiva del Principio de Ambivalencia sobre el estado enriquecido}% ADP102-0200
  \HMTTitleMapEntry{Cristal temporal finito HMT: drive, subarmonía y estabilidad espectral}{Cristal temporal finito HMT: Hamiltoniano periódicamente forzado, respuesta subarmónica y estabilidad espectral}% ADP102-0201
  \HMTTitleMapEntry{Drive construido desde el reloj de 108 estados}{Construcción del forzamiento periódico a partir del sistema temporal de 108 estados}% ADP102-0202
  \HMTTitleMapEntry{Estabilidad espectral finita}{Jet nilpotente de la evolución triádica y prolongación compatible}% ADP102-0203
  \HMTTitleMapEntry{Teorema de confluencia HMT–MD: forma electrogravitatoria, sistema Einstein–Schrödinger y coordinación de acción, materia, campos, geometría e información}{Teorema de confluencia HMT–MD: sistema Einstein–Schrödinger y coordinación electrogravitatoria de acción, materia, campos, geometría e información}% ADP102-0204
  \HMTTitleMapEntry{Independencia horizontal y compatibilidad vertical}{Naturalidad de las realizaciones y conservación vertical de los invariantes del estado enriquecido}% ADP102-0205
  \HMTTitleMapEntry{1.ª confluencia: acoplamiento, vacío y sector electrodébil}{Acoplamiento constitutivo entre el vacío y el sector electrodébil}% ADP102-0206
  \HMTTitleMapEntry{2.ª confluencia: acción, gravedad y forma electrogravitatoria}{Reciprocidad entre acción y gravedad en la forma electrogravitatoria}% ADP102-0207
  \HMTTitleMapEntry{3.ª confluencia: geometría, información y mezcla}{Confluencia entre geometría, información y mezcla}% ADP102-0208
  \HMTTitleMapEntry{4.ª confluencia: pantalla, torsión y memoria helicoidal}{Transducción de la memoria helicoidal en la realización torsional}% ADP102-0209
  \HMTTitleMapEntry{Del determinante electrodébil al operador constitutivo positivo del vacío y su completación espectral \(3\to4\)}{Del determinante electrodébil al operador constitutivo positivo del vacío y su completación espectral 3→4}% ADP102-0210
  \HMTTitleMapEntry{Valor publicado y memoria de la fibra}{Desintegración espectral de la medida de trayectorias y condicionamiento de futuros}% ADP102-0211
  \HMTTitleMapEntry{Fundamentos algebraicos del cosido de fase y de la automedida}{Estado conjunto observador–observable, instrumento reversible y automedida}% ADP102-0212
  \HMTTitleMapEntry{Información par y transporte geométrico impar en la bicapa}{Descomposición par–impar del estado bicapa: información interhoja y transporte geométrico orientado}% ADP102-0214
  \HMTTitleMapEntry{\texorpdfstring{\(3\)}{3} realizaciones de la memoria genealógica: transporte dodecafásico, función de Green no local con generador local y expansión cosmológica}{Realizaciones dinámica, propagatoria y cosmológica de la memoria genealógica: transporte dodecafásico, función de Green y expansión espacial}% ADP102-0215
  \HMTTitleMapEntry{Cadena constitutiva del sector gravitatorio: escala cronológica \texorpdfstring{\(108\)}{108}, reducción tensorial, holonomía y evaluación metrológica de \texorpdfstring{\(G\)}{G}}{Descenso icosaédrico A₅/C₅ y selección del portador gravitatorio}% ADP102-0216
  \HMTTitleMapEntry{Entrelazador con tensores simétricos sin traza y reducción transversal \texorpdfstring{\(12\to5\to5\to2\)}{12→5→5→2}}{Entrelazador con tensores simétricos sin traza y reducción transversal 12→5→5→2}% ADP102-0217
  \HMTTitleMapEntry{Criterios estructurales del selector gravitatorio: transductor dimensional, holonomía enriquecida y separación entre selectores torsional y circular}{Selector holonómico gravitatorio: transducción dimensional, naturalidad, refinamiento y separación de los regímenes torsional y circular}% ADP102-0218
  \HMTTitleMapEntry{Carta decimal declarada y contraste metrológico}{Reconocimiento metrológico posterior de G e independencia causal de la derivación}% ADP102-0219
  \HMTTitleMapEntry{Tétrada, conexión y \(2\) defectos}{Tétrada, conexión de espín y torsión de la realización Cartan–Holst}% ADP102-0220
  \HMTTitleMapEntry{Acción de 1.er orden}{Acción de Palatini–Holst y ecuaciones variacionales}% ADP102-0221
  \HMTTitleMapEntry{Confluencia horizonte–área–información–entropía–acción}{Teorema de confluencia entre energía de horizonte, área, información, entropía y acción}% ADP102-0222
  \HMTTitleMapEntry{\(2\) orientaciones de una misma propagación}{Propagadores avanzado y retardado sobre rutas orientadas}% ADP102-0223
  \HMTTitleMapEntry{Órbita elíptica, acción y \(3\) invariantes J}{Invariantes de la involución temporal y conjugación partícula–antipartícula}% ADP102-0224
  \HMTTitleMapEntry{Cosmología de periodicidad \texorpdfstring{\(108\)}{108}: inversión de hoja, horizonte, Big Bang, rebote torsional y sector oscuro}{Cosmología torsional de periodicidad 108: expansión, retorno con memoria y rebote de Einstein–Cartan}% ADP102-0225
  \HMTTitleMapEntry{Separación tipológica de las \(2\) realizaciones cosmológicas}{Separación entre aceleración de de Sitter y rebote torsional de Einstein–Cartan}% ADP102-0226
  \HMTTitleMapEntry{\(3\) foliaciones exactas de de Sitter}{Foliaciones cosmológicas compatibles con la inversión de hoja}% ADP102-0227
  \HMTTitleMapEntry{Universo con frontera}{Estado de frontera del horizonte cosmológico y memoria de retorno}% ADP102-0228
  \HMTTitleMapEntry{Datos espectrales de una realización supersimétrica: \((\mathcal H_{\rm phys},H,\mathcal Q)\), supercarga y compactificación en dimensión \(11\)}{Realización de teoría M mediante compactificación, cuerdas, branas y dimensión 11}% ADP102-0229
  \HMTTitleMapEntry{Ley General de Masas sobre el atlas \texorpdfstring{\(81/56/13/324\)}{81/56/13/324}: trayectorias, registro, firma, carácter, torres \texorpdfstring{\(90/120\)}{90/120} y sector nulo}{Ley General de Masas sobre el atlas 81/56/13/324: trayectorias, registro, firma, carácter, semigrupo armónico ⟨90,120⟩ y sector nulo}% ADP102-0230
  \HMTTitleMapEntry{Algoritmo universal de realización material e igualador espectral de los lectores \(K_D\) y \(K_\Omega\)}{Construcción functorial de las familias de masa sobre el atlas}% ADP102-0231
  \HMTTitleMapEntry{Bifurcación nula y positividad}{Bifurcación del sector nulo anterior a la rama positiva de masa}% ADP102-0232
  \HMTTitleMapEntry{Portal del dominio exhaustivo}{Construcción exhaustiva del atlas 81/56/13/324 y conservación del sector nulo}% ADP102-0233
  \HMTTitleMapEntry{Dominio interno y problema de realización}{Teorema de exhaustividad del dominio material y estabilidad del sector nulo bajo refinamiento}% ADP102-0234
  \HMTTitleMapEntry{Producto fibrado entre descriptores y rutas}{Producto fibrado entre invariantes de especie y rutas genealógicas}% ADP102-0235
  \HMTTitleMapEntry{Lectores bidireccionales \texorpdfstring{\(K_D\)}{K D} y \texorpdfstring{\(K_\Omega\)}{K Omega}: perfiles, coincidencias, torres \texorpdfstring{\(90/120\)}{90/120} y operadores por multisección}{Lectores bidireccionales \(K_D\) y \(K_\Omega\): perfiles, coincidencias, semigrupo armónico \(\langle90,120\rangle\) y operadores por multisección}% ADP102-0236
  \HMTTitleMapEntry{Lectores \(K_D\) y \(K_\Omega\): perfiles dodecafásicos, dominio de rutas y codominio espectral}{Lectores \(K_D\) y \(K_\Omega\): perfiles dodecafásicos, dominio de rutas y codominio espectral}% ADP102-0237
  \HMTTitleMapEntry{Realizaciones espectrales, escalares y nulas de la Ley General de Masas: familias materiales, transporte CKM--PMNS y sectores fotónico y gluónico}{Realizaciones espectrales, escalares y nulas de la Ley General de Masas: familias materiales, mezcla quark CKM, mezcla neutrínica PMNS y sectores nulos fotónico y gluónico}% ADP102-0238
  \HMTTitleMapEntry{Color, sabores quark y mezcla CKM}{Operadores de masa por fibra, órbita y multisección del atlas}% ADP102-0239
  \HMTTitleMapEntry{Estados compuestos y ley de enlace}{Correspondencia por dominio y codominio entre el inventario material y las realizaciones del atlas}% ADP102-0240
  \HMTTitleMapEntry{Clasificador de especies materiales por fibra, representación, carga, espín, sabor y generación}{Clasificación de las especies materiales por fibra, representación, carga, espín, sabor y generación}% ADP102-0241
  \HMTTitleMapEntry{Tipado probatorio de persistencia, atlas y familias.}{Condiciones de persistencia y correspondencia entre atlas, familias y representaciones}% ADP102-0242
  \HMTTitleMapEntry{Ontología material, carga, espín, estadística, color y sabor}{Definiciones estructurales de partícula, masa, carga, espín, estadística, color, sabor y generación}% ADP102-0243
  \HMTTitleMapEntry{Teorema de cierre tipado de la realización másica}{Teorema de clausura de la realización material mediante secciones persistentes y representaciones finitas}% ADP102-0244
  \HMTTitleMapEntry{Contraste metrológico y archivo exhaustivo de realizaciones}{Reconocimiento metrológico posterior e inventario exhaustivo de las realizaciones materiales}% ADP102-0245
  \HMTTitleMapEntry{Arquitectura natural común de las realizaciones terminales del continuo genealógico}{Teorema de factorización conjunta de las \(5\) realizaciones terminales a través de la estructura discreta del continuo}% ADP102-0246
  \HMTTitleMapEntry{Publicación terminal del continuo único mediante \(5\) lectores no permutables}{Publicación terminal mediante 5 lectores no permutables y factor común del estado enriquecido}% ADP102-0247
  \HMTTitleMapEntry{Clausuras dinámicas del continuo: transporte solenoidal y cociente de calibre}{Dinámica y curvatura del continuo genealógico anteriores a las realizaciones clásicas}% ADP102-0248
  \HMTTitleMapEntry{Dinámica y curvatura del continuo genealógico}{Factorización conjunta de las realizaciones terminales desde el mismo estado del continuo}% ADP102-0249
  \HMTTitleMapEntry{Clausuras algebraicas del continuo y límites de la reconstrucción extensional}{Factorización común, consistencia extensional y límites de la reconstrucción terminal}% ADP102-0250
  \HMTTitleMapEntry{El ejemplo de los ordinales de von Neumann}{Representación ordinal de von Neumann como modelo de consistencia extensional}% ADP102-0251
  \HMTTitleMapEntry{Dominio, mapas y compatibilidades del objeto terminal común}{Vista interna \(G_T^{\mathrm{int}}\) del continuo y aplicación \(\operatorname{App}_{T,d_T}\): conservación de fibra, acarreo, memoria y terminal}% ADP102-0252
  \HMTTitleMapEntry{Despliegue interno de las realizaciones correlativas anterior a su identificación con problemas clásicos}{Emergencia correlativa de las 5 acciones terminales sobre el mismo estado del continuo}% ADP102-0253
  \HMTTitleMapEntry{\(5\) lectores terminales no permutables del continuo genealógico}{Teorema de no permutabilidad de los 5 lectores terminales y clasificación de sus codominios}% ADP102-0254
  \HMTTitleMapEntry{Obstrucciones a la proyección ablativa del continuo: pérdida de fase, orientación o memoria genealógica}{Proyectores positivos de la forma de Guinand–Weil y localización espectral}% ADP102-0255
  \HMTTitleMapEntry{Fijación de la recta crítica por la involución funcional \(s\mapsto1-\bar s\)}{Fijación de la recta crítica por la involución funcional \texorpdfstring{\(s\mapsto 1-\bar{s}\)}{s -> 1 - conjugado(s)}}% ADP102-0256
  \HMTTitleMapEntry{Separación tipada de los ceros de \(\zeta\), \(\xi\) y el determinante espectral}{Separación tipada de los ceros de ζ, ξ y el determinante espectral}% ADP102-0257
  \HMTTitleMapEntry{Representante positivo de la clase residual nula: \(n\equiv0\pmod 9\) y \(\rho_9(n)=9\)}{Representante positivo de la clase residual nula: \(n\equiv0\pmod 9\) y \(\rho_9(n)=9\)}% ADP102-0258
  \HMTTitleMapEntry{Isometría de transporte torcido \(\widetilde U_k\) sobre estados cilíndricos, pesos proyectivos y memoria de Weil}{Isometría de transporte torcido \(\widetilde U_k\) sobre estados cilíndricos, pesos proyectivos y memoria de Weil}% ADP102-0259
  \HMTTitleMapEntry{Invariante graduado del corredor espectral \(4\to2\to6\to54\)}{Invariante graduado del corredor espectral 4→2→6→54}% ADP102-0260
  \HMTTitleMapEntry{Traza acoplada del espectro de Weil y la graduación Moonshine}{Construcción independiente del campo espectral sobre el doble círculo}% ADP102-0261
  \HMTTitleMapEntry{Dominio hidrodinámico y publicación Galerkin}{Dominio solenoidal y conservación de la divergencia nula}% ADP102-0262
  \HMTTitleMapEntry{Dato único de partida y núcleos densos}{Datos iniciales, compatibilidad de frontera y propagación regular}% ADP102-0263
  \HMTTitleMapEntry{Realización finita de la desigualdad de Kalman–Yakubovich–Popov y función auxiliar en la clausura energética}{Desigualdad KYP y control coercivo de la evolución solenoidal}% ADP102-0264
  \HMTTitleMapEntry{Extracción material de la cota de Sobolev}{Estimación de Sobolev y exclusión de la concentración singular}% ADP102-0265
  \HMTTitleMapEntry{La fibra de incidencia es \(1\) coordenada física del estado completo}{Curvatura de fibra y coercividad de la acción de calibre}% ADP102-0266
  \HMTTitleMapEntry{Vacío simple y testigo que alcanza la brecha}{Unicidad del vacío invariante de calibre y saturación de la brecha \(3\kappa_{\mathrm{av}}\) por el autovector \(W_c\)}% ADP102-0267
  \HMTTitleMapEntry{\(4\) reconstrucciones de Osterwalder–Schrader de una dinámica común}{Positividad de Osterwalder–Schrader y reconstrucción cuántica}% ADP102-0268
  \HMTTitleMapEntry{Teoría gauge euclídea y brecha de Yang--Mills}{Brecha espectral del sector de calibre y límite proyectivo}% ADP102-0269
  \HMTTitleMapEntry{Presentación generador a generador de la carta nueva}{Generadores cohomológicos y proyectores compatibles}% ADP102-0270
  \HMTTitleMapEntry{Fibra inicial, Mittag--Leffler y publicación afirmativa}{Lefschetz–Murre y descenso algebraico de las clases de Hodge}% ADP102-0271
  \HMTTitleMapEntry{Bocksteins de todas las páginas}{Homomorfismos de Bockstein de la sucesión espectral: compatibilidad entre páginas y determinación del orden en T del complejo universal}% ADP102-0272
  \HMTTitleMapEntry{Tabla completa Kato--FERS antes del determinante}{Determinante de Kato–FERS, rango y término principal de BSD}% ADP102-0273
  \HMTTitleMapEntry{Separación entre estructura y aplicación}{Vista interna \(G_T\) y lector terminal \(P_T\): independencia lógica entre existencia estructural y realización clásica}% ADP102-0274
  \HMTTitleMapEntry{Obstrucciones comunes a las \(5\) publicaciones terminales: pérdida de dominio, naturalidad, certificado o límite compatible}{Obstrucciones comunes de los lectores terminales: pérdida de dominio, naturalidad o límite compatible}% ADP102-0275
  \HMTTitleMapEntry{Formalización asistida del núcleo algebraico}{Formalización verificable del núcleo algebraico y certificación de identidades finitas}% ADP102-0276
  \HMTTitleMapEntry{Correspondencia entre enunciados, pruebas, certificados y propietarios científicos}{Correspondencia entre enunciados, demostraciones, certificados y fuentes científicas primarias}% ADP102-0277
  \HMTTitleMapEntry{Convenciones de tipado y representación de datos}{Convenciones matemáticas y representación de los datos}% ADP102-0278
  \HMTTitleMapEntry{Heterogeneidad tipada y cobertura del catálogo}{Heterogeneidad de los dominios físicos y exhaustividad del catálogo}% ADP102-0279
  \HMTTitleMapEntry{Fichas individuales de las rutas orientadas}{Despliegue íntegro de las 324 rutas orientadas: correspondencia entre 53 campos canónicos, 57 campos enriquecidos y 49 invariantes publicados por sección}% ADP102-0280
  \HMTTitleMapEntry{Inventario universal tipado de 471 masas y 384 anchuras}{Inventario exhaustivo de 471 masas y 384 anchuras: identificación, procedencia y correspondencia con el atlas}% ADP102-0281
  \HMTTitleMapEntry{Fuente, deduplicación y límites de tipo}{Procedencia de los datos, unicidad de los registros y dominios de validez}% ADP102-0282
  \HMTTitleMapEntry{Código de objeto y lector HMT--MD}{Identificación de objetos y aplicación del lector HMT--MD}% ADP102-0283
  % Correspondencias espirales: sustituciones científicas unívocas aprobadas
  % por cotejo focal. Se modifican sólo los rótulos; los cuerpos permanecen
  % íntegros en sus residencias distribuidas.
  \HMTTitleMapEntry{TPK: selección, transporte, actualización y retorno}{Composición dinámica del Topological Phase Kernel sobre el estado enriquecido: selección trítica, transporte cardinal, actualización de memoria y retorno tipado}% ADP102-0284
  \HMTTitleMapEntry{La transición tipada}{Definición y factorización del operador de transición enriquecida \(U_t\)}% ADP102-0285
  \HMTTitleMapEntry{Composición del acarreo}{Ley de cociclo bimodular del acarreo bajo concatenación de rutas}% ADP102-0286
  \HMTTitleMapEntry{Calendario dodecafásico y no escisión}{Extensión cíclica no escindida \(0\to C_{12}\to C_{108}\to C_9\to0\) y cronología interna del TPK}% ADP102-0287
  \HMTTitleMapEntry{Hélice nonádica de fase y profundidad}{Recubrimiento helicoidal de la fase nonádica y coordenada integral de profundidad de memoria}% ADP102-0288
  \HMTTitleMapEntry{Inmersión geométrica normalizada}{Inmersión \(H_9\) del cociente--residuo nonádico en \(\mathbb R^3\) y cálculo de los invariantes de Frenet}% ADP102-0289
  \HMTTitleMapEntry{Círculos anidados y cilindros}{Secciones circulares y foliación cilíndrica del recubrimiento fase--memoria}% ADP102-0290
  \HMTTitleMapEntry{Hojas orbitales y costura del espacio de suspensión}{Órbitas del producto sesgado y relación de identificación del espacio de suspensión}% ADP102-0291
  \HMTTitleMapEntry{Caso estacionario afín}{Iteración del endomorfismo afín estacionario y clasificación de sus puntos fijos}% ADP102-0292
  \HMTTitleMapEntry{Inversión de rutas}{Antiinvolución del subgrupoide de rutas reversibles y conjugación de la holonomía afín}% ADP102-0293
  \HMTTitleMapEntry{Espejo de constantes y eje fijo}{Involución \(\pi/e\) de la fibra de constantes, invariante \(u^\pi+u^e\) y eje fijo \(\varphi\)}% ADP102-0294
  \HMTTitleMapEntry{La geometría visible como publicación del estado enriquecido}{Evaluación arquimediana del estado enriquecido y defecto de fidelidad genealógica}% ADP102-0295
  \HMTTitleMapEntry{La inversión del orden explicativo}{Factorización causal desde la historia enriquecida hasta la evaluación arquimediana}% ADP102-0296
  \HMTTitleMapEntry{Estado, historia y publicación}{Espacio de estados enriquecidos, categoría de historias TPK y evaluación arquimediana}% ADP102-0297
  \HMTTitleMapEntry{Circunferencia intrínseca y circunferencia proyectada}{Distinción entre órbita radial intrínseca y proyección circular de una trayectoria helicoidal}% ADP102-0298
  \HMTTitleMapEntry{La célula nonádica y el nacimiento bidimensional}{Construcción del soporte nonádico bidimensional \(X_9=(\mathbb Z/9\mathbb Z)^2\)}% ADP102-0299
  \HMTTitleMapEntry{Antecedentes APP de la clasificación radial}{Invariantes APP que inducen el carácter radial acumulado}% ADP102-0300
  \HMTTitleMapEntry{Estado ternario levantado}{Estado TRIT enriquecido por residuo y acarreo: multiplicidad de los representantes del umbral}% ADP102-0301
  \HMTTitleMapEntry{Álgebras cuadráticas}{Familia \(\mathcal A_\tau=\mathbb R[J]/(J^2+\tau I)\) y clasificación elíptica, parabólica e hiperbólica}% ADP102-0302
  \HMTTitleMapEntry{Doble proyección}{Factorización conjunta de las publicaciones arquimediana y excepcional del estado enriquecido}% ADP102-0303
  \HMTTitleMapEntry{Representación local de una transición}{Representación afín \(e\mapsto h_e\) de los generadores de la categoría de rutas}% ADP102-0304
  \HMTTitleMapEntry{Holonomías prefijas}{Productos afines ordenados \texorpdfstring{\(H_j=h_{e_j}\cdots h_{e_1}\)}{H_j = h_ej ... h_e1} y conservación de prefijos}% ADP102-0305
  \HMTTitleMapEntry{Categoría levantada y caracteres elementales}{Categoría de rutas enriquecidas por 2 orientaciones TRIT y definición de los caracteres \(p\) y \(a\)}% ADP102-0306
  \HMTTitleMapEntry{Descomposición balanceada}{Descomposición única \(p=r_3(p)+3c_3(p)\) en residuo balanceado y acarreo}% ADP102-0307
  \HMTTitleMapEntry{Involución}{Acción de \((\tau_+,\tau_-)\mapsto(-\tau_-,-\tau_+)\) sobre los caracteres \(p\) y \(a\)}% ADP102-0308
  \HMTTitleMapEntry{Definición y dominio}{Definición de \(L_p\) sobre \(\mathbb R_{>0}\) y dominio de la inversa \(\exp_p\)}% ADP102-0309
  \HMTTitleMapEntry{Adición deformada}{Ley de composición \(L_p(xy)=L_p(x)+L_p(y)+pL_p(x)L_p(y)\)}% ADP102-0310
  \HMTTitleMapEntry{No exhaustividad y alcance exacto}{Dominio clasificatorio de \(p\in\{2,1,0,-1,-2\}\) y extensión afín general}% ADP102-0311
  \HMTTitleMapEntry{Obstrucción al lector reducido}{Defecto de naturalidad del lector reducido \((p,a,\Lambda,C,m,r,\theta)\)}% ADP102-0312
  \HMTTitleMapEntry{Definición del lector completo}{Lector radial enriquecido \(\mathcal P^{\mathrm{enr}}_{\mathrm{rad},N}\): dominio, codominio y conservación de prefijos}% ADP102-0313
  \HMTTitleMapEntry{Conservación como trazabilidad}{Conservación dinámica del acarreo, la frontera y la memoria por el lector radial}% ADP102-0314
  \HMTTitleMapEntry{Subdivisión cuádruple}{Naturalidad del lector bajo el refinamiento cuádruple \(\operatorname{sd}_4\)}% ADP102-0315
  \HMTTitleMapEntry{Clasificación genealógica y no fidelidad de la publicación}{Clasificación por la firma radial plena y defecto de fidelidad de la evaluación arquimediana}% ADP102-0316
  \HMTTitleMapEntry{Clasificación por 3 incrementos visibles}{Clasificación cinemática mediante los incrementos \((\omega,\beta,\mu)\)}% ADP102-0317
  \HMTTitleMapEntry{3 fibras de una misma evaluación visible}{\(3\) clases genealógicas con una misma imagen arquimediana}% ADP102-0318
  \HMTTitleMapEntry{Ley de tornillo}{Identidad espectral \(V^{-1}\mathscr S V=(2+\sqrt3)\mathrm i\mathscr S\) y ecuación de la suspensión logarítmica}% ADP102-0319
  \HMTTitleMapEntry{Modo trítico común}{Modo singular común de las hojas APP en la carta canónica \(1,\ldots,9\)}% ADP102-0320
  \HMTTitleMapEntry{Ledger material}{Registro genealógico de la ruta material: extensión superviviente, firma entera y carácter adimensional}% ADP102-0321
  \HMTTitleMapEntry{Publicación dimensional posterior}{Realización del operador constitutivo del vacío en los campos \((E_\sigma,B_\sigma)\)}% ADP102-0322
  \HMTTitleMapEntry{Curva plana polar}{Curvatura y longitud de arco de una curva polar \(r=r(\theta)\)}% ADP102-0323
  \HMTTitleMapEntry{Especializaciones}{Evaluación de \(\kappa_{\mathrm F}\) en las familias logarítmica, arquimediana, de Fermat, hiperbólica y lituus}% ADP102-0324
  \HMTTitleMapEntry{Suspensión espacial}{Inmersión espacial \(\Gamma(\theta)\) y criterios diferenciales de Frenet--Lancret}% ADP102-0325
  \HMTTitleMapEntry{Tricromía geométrica}{Coordinación de los regímenes \(J^2=-I,0,+I\) con la geometría de Frenet}% ADP102-0326
  \HMTTitleMapEntry{Conclusión}{Teorema de cierre de las correspondencias espirales genealógicas}% ADP102-0327
  \HMTTitleMapEntry{Semiproducto afín y composición de rutas}{Monoide afín \(\operatorname{End}(U_{\mathrm{APP}})\ltimes U_{\mathrm{APP}}\) y composición ordenada de rutas}% ADP102-0328
  \HMTTitleMapEntry{Tabla exhaustiva de la carta local de 2 orientaciones}{Enumeración exhaustiva de la carta local \((\tau_+,\tau_-)\mapsto(p,a,r_3,c_3)\)}% ADP102-0329
  \HMTTitleMapEntry{5 regiones de \texorpdfstring{\(\pi\)}{pi} y no fidelidad de la evaluación}{Descomposición de la microfibra de \(\pi\) en 5 regiones y multiplicidad de la evaluación arquimediana}% ADP102-0330
  \HMTTitleMapEntry{Álgebra potencia--logarítmica y 5 publicaciones canónicas}{Álgebra de \(L_p\) y clasificación de las 5 familias canónicas \(p\in\{2,1,0,-1,-2\}\)}% ADP102-0331
  \HMTTitleMapEntry{Derivación de las 5 curvas canónicas}{Evaluación de \(\exp_p(u_0+\beta m)\) en las 5 familias canónicas}% ADP102-0332
  \HMTTitleMapEntry{Inversión de hojas}{Conjugación \(p\leftrightarrow-p\) de las hojas potencia--logarítmicas}% ADP102-0333
  \HMTTitleMapEntry{Sector afín general}{Clasificación del sector afín por la pareja \((p,[\Lambda,C])\)}% ADP102-0334
  \HMTTitleMapEntry{Naturalidad, equivariancia y no fidelidad}{Naturalidad bajo truncamiento, equivariancia nonádica y defecto de fidelidad arquimediana}% ADP102-0335
  \HMTTitleMapEntry{El lector de prefijos}{Lector de prefijos afines y conservación del registro enriquecido}% ADP102-0336
  \HMTTitleMapEntry{Refinamientos no triviales}{Criterio de compatibilidad para refinamientos generales de rutas}% ADP102-0337
  \HMTTitleMapEntry{Equivarianza, no invariancia}{Equivarianza del cociclo radial bajo la holonomía \(\Gamma_9\)}% ADP102-0338
  \HMTTitleMapEntry{3 demostraciones de no fidelidad}{\(3\) obstrucciones explícitas a la fidelidad de la evaluación arquimediana}% ADP102-0339
  \HMTTitleMapEntry{Mismo extremo, distinto orden afín}{Coincidencia del producto final y dependencia del orden de los factores afines}% ADP102-0340
  \HMTTitleMapEntry{Misma curva, distinta región genealógica}{Coincidencia de la imagen arquimediana entre regiones genealógicas distintas}% ADP102-0341
  \HMTTitleMapEntry{Clasificación por nivel matemático}{Estratificación por dominio: historia, lector, evaluación arquimediana y realización física}% ADP102-0342
  \HMTTitleMapEntry{Clasificación por giro, radio y memoria}{Clasificación por los incrementos \((\omega,\beta,\mu)\)}% ADP102-0343
  \HMTTitleMapEntry{Realizaciones posteriores y controles de tipo}{Realizaciones electromagnética y material con separación de dominios y codominios}% ADP102-0344
  \HMTTitleMapEntry{Centro electrónico y lector posterior}{Compatibilidad del lector radial con el punto fijo electrónico \((5,5)\) y el atlas \(81/56/13/324\)}% ADP102-0345
  \HMTTitleMapEntry{Matriz cerrada de controles negativos}{Matriz de condiciones de validez por nivel causal}% ADP102-0346
  \HMTTitleMapEntry{Primera prolongación y frontera coinductiva}{Prolongación coinductiva inicial desde \(w_6\) hasta la frontera de compatibilidad de orden \(1\)}% ADP102-0347
  \HMTTitleMapEntry{Vía B: núcleo de vacancias \texorpdfstring{\(E_{21/22}\)}{E21/22}}{Raíz simple \(\xi_9\) del núcleo analítico de orden 9 y comparación exacta \(\operatorname{Tr}_9(\xi_9^{-1})=\operatorname{Tr}_9(\alpha_{\mathrm{HMT}}^{-1})\)}% ADP102-0348
  % Purismo residual: rótulos activos localizados en el índice del sucesor.
  % Las sustituciones siguientes no alteran el nivel jerárquico ni el cuerpo.
  \HMTTitleMapEntry{Cronología observable de los dos cursores}{Evolución cronológica del par de cursores y determinación del estado observable}% ADP102-0349
  \HMTTitleMapEntry{Emisor finito de dos cursores}{Emisor finito asociado al par de cursores y generación de estados admisibles}% ADP102-0350
  \HMTTitleMapEntry{De una célula a dos coordenadas: nacimiento de APP}{Construcción de APP mediante la descomposición de 1 célula en 2 coordenadas}% ADP102-0351
  \HMTTitleMapEntry{La rama de \(20\) bloques}{Prolongación de la sección superviviente mediante 20 bloques compatibles}% ADP102-0352
  \HMTTitleMapEntry{Del prefijo a un punto real}{Evaluación arquimediana del prefijo compatible y determinación del punto real}% ADP102-0353
  \HMTTitleMapEntry{Una emisión reproducible de doce tríadas}{Emisión reproducible de 12 tríadas y reconstrucción de su estado enriquecido}% ADP102-0354
  \HMTTitleMapEntry{Construcción del carácter de clausura y trayectorias genealógicas estructural de \(\pi\)}{Construcción del carácter de clausura y de las trayectorias genealógicas de \(\pi\)}% ADP102-0355
  \HMTTitleMapEntry{La pendiente 1/5 forzada por las \(5\) regiones}{Determinación de la pendiente 1/5 mediante el balance de las 5 regiones orientadas}% ADP102-0356
  \HMTTitleMapEntry{Ventana dodecafásica del estado firmado y clausura reversible del ciclo}{Restricción finita del estado dodecafásico firmado y clausura reversible del ciclo}% ADP102-0357
  \HMTTitleMapEntry{Ventana dodecafásica finita}{Dominio finito de la clausura dodecafásica}% ADP102-0358
  \HMTTitleMapEntry{Prolongación remota \texorpdfstring{\(\mathsf T_{\alpha,\infty}\)}{T alpha infinito} de la vía entera}{Prolongación coinductiva \(\mathsf T_{\alpha,\infty}\) de la clausura dodecafásica entera}% ADP102-0359
  \HMTTitleMapEntry{Teorema de los cuatro canales enteros}{Teorema de clasificación de los 4 canales enteros}% ADP102-0360
  \HMTTitleMapEntry{Primer codificador intrínseco y su límite}{Codificador intrínseco inicial y existencia de su límite}% ADP102-0361
  \HMTTitleMapEntry{Hipótesis de tipado y controles negativos}{Condiciones de tipado y criterios de incompatibilidad de la teoría de vacancias nonádicas}% ADP102-0362
  \HMTTitleMapEntry{Vía cronológica de longitud, acción y espectro}{Composición cronológica de longitud, acción y espectro en la determinación de \(G\)}% ADP102-0363
  \HMTTitleMapEntry{Operador torcido y control negativo determinista}{Operador con torsión y obstrucción determinista a la continuación multiplicativa}% ADP102-0364
  \HMTTitleMapEntry{Estratificación matemática del objeto Majorana en cuatro niveles}{Estratificación matemática del objeto Majorana en 4 niveles}% ADP102-0365
  \HMTTitleMapEntry{Seis caracteres y cuatro torsiones de trenza}{Clasificación de 6 caracteres por 4 clases de torsión de trenza}% ADP102-0366
  \HMTTitleMapEntry{obstrucciones y condiciones de no degeneración específicos}{Obstrucciones específicas y condiciones de no degeneración}% ADP102-0367
  \HMTTitleMapEntry{Quark arriba y antiquark arriba}{Sectores quark \(u\) y antiquark \(\bar u\)}% ADP102-0368
  \HMTTitleMapEntry{Quark abajo y antiquark abajo}{Sectores quark \(d\) y antiquark \(\bar d\)}% ADP102-0369
  \HMTTitleMapEntry{Controles de refutación.}{Condiciones de validez e incompatibilidad de la realización equivariante de especies}% ADP102-0370
  \HMTTitleMapEntry{El cuadrado local de \texorpdfstring{\(e\)}{e} y la ruptura de memoria}{Cuadrado posicional local de \(e\) y diferenciación genealógica de estados con igual bloque visible}% ADP102-0371
  \HMTTitleMapEntry{La ley general antes de sus realizaciones por especie}{Antecedencia causal de la Ley General de Masas respecto de las realizaciones por especie}% ADP102-0372
  \HMTTitleMapEntry{Fundamento unificado de la masa}{Factorización extensión–ruta y conservación del registro genealógico}% ADP102-0373
  \HMTTitleMapEntry{La banda de ruta}{Definición de la banda de ruta mediante la órbita de una sección persistente}% ADP102-0374
  \HMTTitleMapEntry{El canal de incidencia es curvatura local, no un factor auxiliar}{Interpretación del canal de incidencia como curvatura local del sistema de calibre}% ADP102-0375
  % Publicaciones terminales del capítulo 95: se conserva cada desarrollo
  % completo y se sustituye la enumeración narrativa por objeto, operación y resultado.
  \HMTTitleMapEntry{1.ª resolución: Riemann y la positividad como complemento construido}{Publicación espectral–polar del continuo mediante el complemento positivo de Guinand–Weil}% ADP102-0376
  \HMTTitleMapEntry{Existencia, regularidad y unicidad global para Navier–Stokes mediante cota uniforme y compacidad}{Publicación solenoidal del continuo: cota coerciva, compacidad y regularidad global}% ADP102-0377
  \HMTTitleMapEntry{3.ª resolución: Yang–Mills y la brecha del sector de curvatura}{Publicación de calibre del continuo: positividad por reflexión y brecha del sector de curvatura}% ADP102-0378
  \HMTTitleMapEntry{4.ª resolución: Hodge y la construcción efectiva de la preimagen algebraica}{Publicación cohomológica del continuo: levantamiento compatible y construcción efectiva de la preimagen algebraica}% ADP102-0379
  \HMTTitleMapEntry{5.ª resolución: BSD y la coincidencia de las publicaciones analítica y aritmética}{Publicación determinantal del continuo: comparación Kato–Poitou–Tate e identidad del término principal de Birch–Swinnerton–Dyer}% ADP102-0380
  \HMTTitleMapEntry{La identidad de scattering que conserva la potencia cruzada}{Identidad de dispersión y conservación de la potencia cruzada}% ADP102-0381
  \HMTTitleMapEntry{Persistencia del objeto terminal bajo el paso al límite}{Persistencia del estado terminal solenoidal bajo el paso al límite}% ADP102-0382
  \HMTTitleMapEntry{Existencia, regularidad y unicidad global de la solución de Navier–Stokes obtenida por cota uniforme y compacidad}{Teorema de clausura solenoidal global mediante cota uniforme y compacidad}% ADP102-0383
  \HMTTitleMapEntry{El operador universal de extensión}{Operador universal de extensión de clases cohomológicas compatibles}% ADP102-0384
  \HMTTitleMapEntry{Factorización de los entrelazadores terminales RH–NS–YM–Hodge–BSD a través del continuo genealógico común}{Pérdidas de información específicas de las 5 publicaciones clásicas y conservación en el estado enriquecido}% ADP102-0385
  \HMTTitleMapEntry{Factorización de RH–NS–YM–Hodge–BSD como publicaciones terminales del continuo genealógico común}{Teorema de factorización terminal de RH–NS–YM–Hodge–BSD a través del continuo genealógico común}% ADP102-0386
  \HMTTitleMapEntry{Emergencia correlativa de \texorpdfstring{\(5\)}{5} acciones intrínsecas no exhaustivas}{Descomposición expositiva de la acción correlativa en 5 vistas intrínsecas del mismo estado}% ADP102-0387
  \HMTTitleMapEntry{Una misma frontera: flujo interior y tensión gauge}{Factorización común del operador de frontera en las realizaciones solenoidal y de calibre}% ADP102-0388
  \HMTTitleMapEntry{El estado común de frontera}{Estado enriquecido de frontera y descomposición paritaria del transporte}% ADP102-0389
  \HMTTitleMapEntry{El operador mínimo y su espectro}{Complejo de borde mínimo y equivalencia espectral de los laplacianos de grados 0 y 1}% ADP102-0390
  \HMTTitleMapEntry{Lectura hidrodinámica}{Realización solenoidal del operador de frontera: contracción del flujo y conservación de energía}% ADP102-0391
  \HMTTitleMapEntry{Lectura Yang--Mills}{Realización de calibre del operador de frontera: coercividad, cociente gauge y brecha espectral}% ADP102-0392
  % Promoción jerárquica de los certificados causales de \(\Gamma_9\) y de
  % la generación trítica de los canales enteros.
  \HMTTitleMapEntry{Filtración monodrómica y cálculo a profundidad prescrita}{Certificación finita de Γ₉: 396 niveles, 2.376 trits, 43 vueltas completas, 387 pares \((K,K+9)\) por canal y fibra ambiente de 729 elementos}% ADP102-0393
  \HMTTitleMapEntry{Censo finito de transiciones y generación decimal}{Censo exhaustivo de 396 transiciones, 2.376 trits, 43 vueltas y 387 pares de igual fase}% ADP102-0394
  \HMTTitleMapEntry{Estructura entera de los lectores de constantes}{Clasificación aritmética de los canales enteros y compatibilidad con la semicorona trítica}% ADP102-0395
}
\let\HMTNativeChapter\chapter
\let\HMTNativeSection\section
\let\HMTNativeSubsection\subsection
\let\HMTNativeSubsubsection\subsubsection
\let\HMTNativeParagraph\paragraph
\let\HMTNativeSubparagraph\subparagraph
\RenewDocumentCommand{\chapter}{s o m}{%
  \HMTResolvePublicTitle{#3}%
  \IfBooleanTF{#1}{\HMTNativeChapter*{\HMTResolvedTitle}}{%
    \IfNoValueTF{#2}{\HMTNativeChapter{\HMTResolvedTitle}}{\HMTNativeChapter[\HMTResolvedTitle]{\HMTResolvedTitle}}}}
\RenewDocumentCommand{\section}{s o m}{%
  \HMTResolvePublicTitle{#3}%
  \IfBooleanTF{#1}{\HMTNativeSection*{\HMTResolvedTitle}}{%
    \IfNoValueTF{#2}{\HMTNativeSection{\HMTResolvedTitle}}{\HMTNativeSection[\HMTResolvedTitle]{\HMTResolvedTitle}}}}
\RenewDocumentCommand{\subsection}{s o m}{%
  \HMTResolvePublicTitle{#3}%
  \IfBooleanTF{#1}{\HMTNativeSubsection*{\HMTResolvedTitle}}{%
    \IfNoValueTF{#2}{\HMTNativeSubsection{\HMTResolvedTitle}}{\HMTNativeSubsection[\HMTResolvedTitle]{\HMTResolvedTitle}}}}
\RenewDocumentCommand{\subsubsection}{s o m}{%
  \HMTResolvePublicTitle{#3}%
  \IfBooleanTF{#1}{\HMTNativeSubsubsection*{\HMTResolvedTitle}}{%
    \IfNoValueTF{#2}{\HMTNativeSubsubsection{\HMTResolvedTitle}}{\HMTNativeSubsubsection[\HMTResolvedTitle]{\HMTResolvedTitle}}}}
\RenewDocumentCommand{\paragraph}{s o m}{%
  \HMTResolvePublicTitle{#3}%
  \IfBooleanTF{#1}{\HMTNativeParagraph*{\HMTResolvedTitle}}{%
    \IfNoValueTF{#2}{\HMTNativeParagraph{\HMTResolvedTitle}}{\HMTNativeParagraph[\HMTResolvedTitle]{\HMTResolvedTitle}}}}
\RenewDocumentCommand{\subparagraph}{s o m}{%
  \HMTResolvePublicTitle{#3}%
  \IfBooleanTF{#1}{\HMTNativeSubparagraph*{\HMTResolvedTitle}}{%
    \IfNoValueTF{#2}{\HMTNativeSubparagraph{\HMTResolvedTitle}}{\HMTNativeSubparagraph[\HMTResolvedTitle]{\HMTResolvedTitle}}}}
% Las utilidades de degradación posteriores deben conservar el mapa de títulos.
\let\HMTBookChapter\chapter
\let\HMTBookSection\section
\let\HMTBookSubsection\subsection
\let\HMTBookSubsubsection\subsubsection
\let\HMTBookParagraph\paragraph
\let\HMTBookSubparagraph\subparagraph
\makeatother
